% TOP_PROBLEM: {"problemId": "problem.the-atiyah-conjecture-on-the-integrality-of-l2-betti-numbers", "canonicalTitle": "The Atiyah Conjecture on the Integrality of L2-Betti Numbers", "releaseRank": 334, "primaryDomain": "algebra_representation_category", "primaryDomainLabel": "Algebra, representation & category theory", "displayStatus": "open", "releaseStatus": "open"}
% !TeX program = lualatex
% Definition and sources only; this file is not an LLM solution attempt.
\documentclass[11pt]{article}
\usepackage[margin=25mm]{geometry}
\usepackage{fontspec}
\setmainfont{Latin Modern Roman}
\usepackage{amsmath,amssymb,mathtools}
\usepackage{unicode-math}
\setmathfont{Latin Modern Math}
\usepackage{xurl,hyperref}
\hypersetup{colorlinks=true,urlcolor=blue}
\setcounter{secnumdepth}{0}
\setlength{\parindent}{0pt}
\setlength{\parskip}{5pt}
\setlength{\emergencystretch}{3em}
\begin{document}
\section*{\texorpdfstring{The Atiyah Conjecture on the Integrality of L2-Betti Numbers}{The Atiyah Conjecture on the Integrality of L2-Betti Numbers}}\label{334}
\subsection{Definitions and mathematical statement}
Let $X$ be a finite CW-complex and $G=\pi_1(X)$ torsion-free. Lifting cells to its universal cover gives a Hilbert chain complex with spaces $C_i^{(2)}(\widetilde X)\cong\ell^2(G)^{c_i}$ and bounded cellular boundary maps $d_i$. Its reduced $L^2$ homology is
\[
 \mathcal H_i^{(2)}=\ker d_i\big/\overline{\operatorname{im}d_{i+1}},
 \qquad b_i^{(2)}(\widetilde X)=\dim_{\mathcal N(G)}\mathcal H_i^{(2)}.
\]
The von Neumann dimension is computed using the canonical group trace on the orthogonal projection representing this Hilbert module, not an ordinary vector-space dimension. The selected Atiyah assertion is $b_i^{(2)}(\widetilde X)\in{\mathbb{Z}}_{\ge0}$ for every $i$ and every such $X$. No assertion about groups with torsion or denominator bounds for their invariants is substituted.
\par\smallskip\textit{Formulation and concept references:} \href{https://arxiv.org/abs/math/0001101}{[S1]}.

\subsection{Short English statement}
For every finite cell complex with torsion-free fundamental group, must all von Neumann dimensions of its square-integrable homology be integers?
\subsection{Sources}
\begin{itemize}
\item[S1]T. Schick, Math. Ann. 317 (2000), 727-750.
\url{https://arxiv.org/abs/math/0001101}.
\textit{Formulation source.}
\end{itemize}
\end{document}
